\subsection{Calculating Indirect Genetic Support From Gene Annotations and Genetic Evidence}
\begin{figure}[h!]
    \centering
    \includegraphics[width=.85\linewidth]{figures/indirect-support/real/PIGEAN-Figure-1.pdf}
    \caption{Conceptual overview of the PIGEAN framework. (A) Demostration of top-down prioritization methods from GWAS, WES and/or Mendelian disease databases. (B) Demostration of bottom-up priortization methods from gene annotations. (C) Graphical Model of PIGEAN with a depiction of fusing direct and indirect genetic evidence. (D) Positioning each gene across the continum of direct and indirect support. (E) Leveraging this continuum to power downstream analysis pipelines which can be leveraged in agentic tooling (F) to yeild AI-driven analysis at scale.}
    \label{fig:pigean-concept}
\end{figure}
To map gene-trait relationships across the entire phenome and genome, we sought to infer a probability of relevance for any gene  to any trait. Doing so using available human genetic associations and non-human-genetic data posed several challenges. First, we needed to infer relationships across thousands of traits accurately, robustly, and efficiently. Second, the traits are heterogeneous in terms of genetic association data available (e.g. common variant associations for GWAS, Mendelian gene-disease validity assertions for rare diseases) and genetic architecture. Third, currently significant gene-trait associations cover only a small fraction of trait-relevant genes. Fourth, non-human-genetic datasets (i.e. curated pathway databases, results from animal model experiments, text-mined relationships, or gene expression enrichment analyses) are large in number and of varying quality, and it is unclear which datasets are useful for predicting trait-relevant genes. Finally, non-human-genetic datasets are biased toward well-studied genes and well-studied biological mechanisms.

To address these challenges, we developed a statistical model of a quantity we term a gene's ``indirect genetic support'', and a method termed PIGEAN (Priors Inferred from GEne ANnotations) for homogeneously inferring across common and rare traits. Indirect genetic support combines the unbiased rigor of genetic associations with the breadth and scope of gene annotations, which encode non-genetic-data either as gene sets (binary annotations) or gene quantitative vectors (continuous annotations) and can be constructed in a high-throughput manner from many biomedical data sources (\ref{sec:methods}). Indirect genetic support is an emergent property of a simple statistical model of observed genetic associations, in which gene annotations affect the probability that a gene is relevant to a trait, with effect sizes unknown and different for each annotation and each trait, and in which trait-relevant genes are more likely to produce genetic associations than are non-trait-relevant genes (\ref{sec:methods}).  This model statistically generalizes several widely used heuristic approaches for prioritizing genes from genetic association studies (\ref{sec:methods}), and it enables us to demonstrate mathematically that genetic associations not only directly implicate the genes for which they are observed but also indirectly implicate genes with shared annotations. PIGEAN applies Bayesian inference to this model to jointly infer gene annotation effect sizes and gene trait-relevance probabilities given a library of gene annotations and either GWAS summary statistics or rare disease gene lists, in so doing (a) using human genetic associations to select annotations with non-zero effects and calibrate their effect sizes and (b) using genetic annotations to identify additional trait-relevant genes beyond those with direct association evidence. Priors on the latent parameters enforce strict regularization that control biases towards well-studied genes or profligate annotations, as well as enforcing model sparsity and interpretability (\ref{sec:methods}).

We studied PIGEAN across a universe of 5,030 complex traits and 1,526 rare disease groups. For complex traits, we integrated full summary statistic files from XXX GWAS covering XXX traits, as well as curated association lists from XXX GWAS across XXX traits. We used an extension of a previously published approach \cite{dornbos_evaluating_2022} to calculate each gene's direct genetic support from variant-level associations (\ref{sec:methods}).  For rare diseases, we mapped curated disease–gene confidence levels from Orphanet to probabilities of disease relevance (\ref{sec:methods}).  

To evaluate the contribution of non-genetic evidence, we constructed multiple PIGEAN models, each defined by a specific set of geneset libraries available during inference. These libraries collectively included 144,627 gene annotations spanning diverse sources (\ref{Supplementary Table 2}): knockout mouse phenotypes (11,502 Mammalian Phenotype Ontology terms), membership in pathways and curated gene sets (29,003 from the Molecular Signatures Database), co-occurrence in PubMedCentral manuscript figures (19,391 sets from the Pathway Figure OCR project), co-occurrence with terms in PubMed abstracts (25,126 MeSH terms), network interaction partners (37,156 STRING neighbors), and single-cell gene expression patterns (41,840 experimental clusters curated as part of the PoPS method).  

The resulting map (Figure \ref{fig:full_map}), generated from the top-performing model identified via model selection (See \textit{Validation of Resource}), produced a complete numerical matrix of gene–trait relevance probabilities. This included 135,188 gene-trait relationships (21 per trait and 7 per gene) with a posterior probability (PP) above 90\% (i.e. false discovery rate < 0.1), 295,547 gene–trait relationships (45 per trait and 16 per gene) with PP>50\%, and 10.49M significant (\ref{sec:methods}) gene set–trait associations. Each gene-trait relationship is supported by provenance that includes the specific SNP associations or curated gene records  that produce direct support, and the specific gene set-trait relationsips that produce indirect support. The full map is available for download, for interactive browsing via a web interface, or through AI-driven query to filter and explore results in response to high-level questions.  